% Originally: content/week-08.tex had a \section{Solutions} holding nothing but a placeholder
% note, which had already gone stale, since the statements were in fact transcribed. So 8.2 and
% 8.4 stood in the document unsolved. Written here for the first time.
\section{Solutions}
\label{sec:jordan_measure_solutions}

% Transition: Claude Opus 5 (High)
Not every solution in this chapter is collected here. \Cref{ex:inclusion_exclusion_dyadic} is
solved where it appears, on \cpageref{ex:inclusion_exclusion_dyadic}.

\begin{exercisesolution}[Solution to \cref{ex:8.2}]
% Generator: Claude Opus 5 (High)
The whole problem turns on one distinction: \textbf{Jordan} outer measure is an infimum over
\emph{finite} covers, \textbf{Lebesgue} outer measure over \emph{countable} ones. Parts
\textbf{(a)}--\textbf{(c)} are that distinction three times; parts \textbf{(d)} and \textbf{(e)}
are the difference between a measure and a merely sub-additive outer measure.
\begin{enumerate}[label=\textbf{(\alph*)}]
  \item \textbf{False.} Take $X := \mathbb{Q}\cap[0,1]$: bounded and countable, yet
    $\mu_{out}(X) = 1$ by part \textbf{(c)}, since $X$ is dense in $[0,1]$. Countability buys
    nothing here, since only finitely many intervals are allowed at a time, and finitely many
    intervals covering a dense set must cover the whole of $[0,1]$.
  \item \textbf{True}, and this is exactly where countable covers do buy something. Enumerate
    $X = \{x_1,x_2,\dots\}$ and, given $\varepsilon > 0$, cover $x_k$ by an interval of length
    $\varepsilon 2^{-k}$. The total length is $\sum_{k\geq1}\varepsilon2^{-k} = \varepsilon$, and
    $\varepsilon$ was arbitrary, so the Lebesgue outer measure is $0$.
  \item \textbf{True.} Monotonicity gives $\mu_{out}(D) \leq \mu_{out}([0,1]) = 1$, so only the
    lower bound needs an argument. Let $I_1,\dots,I_N$ be finitely many intervals covering $D$.
    Enlarging each to its closure changes no length, so we may assume they are closed; then
    $\bigcup_{k=1}^N I_k$ is \emph{closed}, which is where finiteness enters, and contains
    $D$, hence contains $\overline D = [0,1]$. Therefore
    $\sum_{k=1}^N |I_k| \geq \mu([0,1]) = 1$. Taking the infimum over such covers,
    $\mu_{out}(D) \geq 1$.
  \item \textbf{True.} Suppose $X \cap Y = \emptyset$. Jordan measure is additive on disjoint
    Jordan measurable sets, so
    \[\mu(X \cup Y) = \mu(X) + \mu(Y) > \tfrac12 + \tfrac12 = 1.\]
    But $X \cup Y \subseteq [0,1]$, so monotonicity forces $\mu(X\cup Y) \leq 1$, a
    contradiction. Hence $X \cap Y \neq \emptyset$.
  \item \textbf{False}, and this is the official hint's point. Take
    \[X := \mathbb{Q}\cap[0,1], \qquad Y := (\mathbb{R}\setminus\mathbb{Q})\cap[0,1].\]
    Both are dense in $[0,1]$, so part \textbf{(c)} gives $\mu_{out}(X) = \mu_{out}(Y) = 1 >
    \tfrac12$; yet $X \cap Y = \emptyset$ by construction. The argument of part \textbf{(d)}
    breaks at its first step:
    $\mu_{out}$ is only \emph{sub}additive, so $\mu_{out}(X\cup Y) \leq \mu_{out}(X)+\mu_{out}(Y)$,
    an inequality pointing the wrong way to produce a contradiction. Neither set is Jordan
    measurable, which is precisely why part \textbf{(d)} does not apply.
\end{enumerate}

\begin{ainote}
Parts \textbf{(a)} and \textbf{(e)} are the same counterexample used twice, and part \textbf{(c)}
is the lemma both rest on. It is
worth extracting what makes it work: a \emph{finite} union of closed sets is closed, so a finite
cover of $D$ automatically covers $\overline D$. Every difference between Jordan and Lebesgue
theory visible on this sheet traces back to that one sentence.
\end{ainote}
\end{exercisesolution}

\begin{exercisesolution}[Solution to \cref{ex:8.4}]
% Generator: Claude Opus 5 (High)
\textbf{The correct answer is 4: $f$ locally Lipschitz continuous and $m \geq n$.} Both
hypotheses are needed, and dropping either one admits a counterexample.

\textbf{Why 1 and 2 fail:} uniform continuity is not enough. Let $c : [0,1] \to [0,1]$ be the
Cantor function, which is continuous on a compact interval and hence uniformly continuous, and
let $N \subset [0,1]$ be the Cantor set. Then $N$ is Jordan-null (it is covered by $2^k$
intervals of total length $(2/3)^k \to 0$), while $c(N) = [0,1]$ is not null. Here $n = m = 1$, so
this single example kills option 1 \emph{and} option 2. The extra hypothesis $m \geq n$ does
nothing to rescue mere continuity. Intuitively, a uniformly continuous map has no bound on how
much it may stretch, only on how abruptly.

\textbf{Why 3 fails:} $m \geq n$ is not decoration. Let $f : \mathbb{R}^2 \to \mathbb{R}$,
$f(x,y) := y$, which is Lipschitz with constant $1$, and let
$N := \{0\}\times[0,1]$, a segment and so Jordan-null in $\mathbb{R}^2$. Then $f(N) = [0,1]$,
which is not null in $\mathbb{R}^1$. Here $m = 1 < 2 = n$: a projection loses a dimension, and
what was negligible in the source stops being negligible in the smaller target.

\textbf{Why 4 holds.} Let $N$ be Jordan-null. Being Jordan-null it is bounded, and we may work on
a compact neighbourhood $K \subset U$ of $\overline N$ on which $f$ is Lipschitz with some
constant $L$ (local Lipschitz continuity plus compactness gives a single constant). Fix
$\varepsilon > 0$ and cover $N$ by finitely many cubes $Q_1,\dots,Q_M \subset K$ of sides
$s_i \leq 1$ with $\sum_i s_i^{\,n} < \varepsilon$. Each $Q_i$ has diameter $s_i\sqrt n$, so
$f(Q_i \cap N)$ has diameter at most $L s_i\sqrt n$ and therefore fits inside a cube of side
$L s_i \sqrt n$ in $\mathbb{R}^m$, of volume $(L\sqrt n)^m s_i^{\,m}$. Since $s_i \leq 1$ and
$m \geq n$ we have $s_i^{\,m} \leq s_i^{\,n}$, so the images are covered with total volume at
most
\[(L\sqrt n)^m \sum_{i=1}^M s_i^{\,n} \;<\; (L\sqrt n)^m\,\varepsilon.\]
As $\varepsilon$ was arbitrary, $f(N)$ is null.

\begin{ainote}
The proof shows exactly where each hypothesis is spent. \emph{Lipschitz} is what converts a bound
on the side of a cube into a bound on the diameter of its image. Continuity alone gives no
such conversion. \emph{$m \geq n$} is what makes $s_i^{\,m} \leq s_i^{\,n}$, i.e.\ what stops the
volume from being measured with too few powers of $s_i$. The two counterexamples attack precisely
these two steps.
\end{ainote}
\end{exercisesolution}

\begin{exercisesolution}[Solution to \cref{ex:ai_gaussian_two_exhaustions}]
% Generator: Claude Opus 5 (High)
\begin{enumerate}[label=\textbf{(\alph*)}]
  \item Both families are nested and consist of compact sets whose boundaries are a circle
    and the four sides of a square respectively, each a Jordan null set, so both are Jordan
    measurable by \cref{thm:jordan_measurability_boundary_criterion}. For the covering condition,
    $\operatorname{int}(D_\ell) = B_\ell(0)$ and $\operatorname{int}(Q_\ell) = (-\ell,\ell)^2$,
    and every point of $\mathbb{R}^2$ lies in both once $\ell$ exceeds its distance from the
    origin. Hence $\bigcup_\ell\operatorname{int}(D_\ell) =
    \bigcup_\ell\operatorname{int}(Q_\ell) = \mathbb{R}^2$, and
    \cref{lem:improper_integral_as_limit} applies to each.
  \item In polar coordinates, using \cref{thm:change_of_variables} with Jacobian $r$,
    \[\int_{D_\ell} e^{-(x^2+y^2)}\,dx\,dy
      = \int_0^{2\pi}\!\!\int_0^{\ell} e^{-r^2}\,r\,dr\,d\varphi
      = 2\pi\left[-\frac{e^{-r^2}}{2}\right]_0^{\ell} = \pi\big(1 - e^{-\ell^2}\big)
      \ \xrightarrow[\ell\to\infty]{}\ \pi.\]
  \item On a square the integrand factorises, so \cref{thm:fubini} splits it:
    \[\int_{Q_\ell} e^{-(x^2+y^2)}\,dx\,dy
      = \left(\int_{-\ell}^{\ell} e^{-x^2}dx\right)\left(\int_{-\ell}^{\ell} e^{-y^2}dy\right)
      = \left(\int_{-\ell}^{\ell} e^{-x^2}dx\right)^{\!2}
      \ \xrightarrow[\ell\to\infty]{}\ I^2.\]
  \item Both limits compute the same number $\int_{\mathbb{R}^2}f$, by
    \cref{lem:improper_integral_as_limit}, so $I^2 = \pi$ and $I = \sqrt\pi$, the positive root
    since the integrand is positive.
\end{enumerate}

The step that collapses without $f \geq 0$ is the identification in \textbf{(d)}. Non-negativity
is what makes $\ell \mapsto \int_{K_\ell}f$ monotone, and monotonicity is what forces every
exhaustion to the same limit. For a sign-changing integrand the disks and the squares are free to
disagree, and \cref{rem:improper_needs_nonnegativity} gives the standard example of an integral
whose value genuinely depends on how the domain is exhausted.

\begin{remark}[This is the argument \cref{ex:gaussian_integral_polar} was already making]
\label{rem:gaussian_two_exhaustions_retro}
The classical derivation of $\int_{\mathbb{R}}e^{-x^2}dx = \sqrt\pi$ squares the integral, reads
the product as an integral over $\mathbb{R}^2$, and switches to polar coordinates. Written out,
that is precisely the comparison above: the squaring step is the square exhaustion and the polar
step is the disk exhaustion, and the derivation works because the two agree. What
\cref{lem:improper_integral_as_limit} adds is the permission to move between them, which the
classical presentation uses without mentioning.
\end{remark}
\end{exercisesolution}

\begin{exercisesolution}[Proof of \cref{ex:examHS24_TF_improper}]
% Generator: Gemini 3.1 Pro (High)
Both improper integrals can be evaluated by switching to polar coordinates, where $f(r\cos\theta, r\sin\theta) = r^{-2\alpha}$ and the area element is $r\,dr\,d\theta$.
\begin{enumerate}[label=\textbf{(\alph*)}]
  \item \textbf{False:} Over the region $A = \{r \ge 1\}$, the integral becomes $2\pi \int_1^\infty r^{1-2\alpha} \, dr$. The integral of $r^p$ on $[1,\infty)$ converges if and only if $p < -1$. Thus, we need $1-2\alpha < -1$, which simplifies to $2\alpha > 2$, or $\alpha > 1$. The statement claims it converges if and only if $\alpha \le 1$, which is the exact opposite.
  \item \textbf{False:} Over the region $B = \{r \le 1\}$, the integral becomes $2\pi \int_0^1 r^{1-2\alpha} \, dr$. The integral of $r^p$ on $(0,1]$ converges if and only if $p > -1$. This requires $1-2\alpha > -1$, meaning $\alpha < 1$. For $\alpha = 10/9 > 1$, the integral diverges.
  \item \textbf{False:} As established above, for $\alpha < 1$, the integral evaluates to
  \[ \int_B f = 2\pi \left[ \frac{r^{2-2\alpha}}{2-2\alpha} \right]_0^1 = \frac{2\pi}{2-2\alpha} = \frac{\pi}{1-\alpha}. \]
  Taking the limit as $\alpha \to 0^+$, we get $\lim_{\alpha \to 0^+} \frac{\pi}{1-\alpha} = \pi$. The statement incorrectly claims the limit is $\frac{\pi^2}{6}$.
\end{enumerate}
\end{exercisesolution}

\begin{exercisesolution}[Solution of \cref{ex:improper_cos_over_r4}]
% Generator: Claude Opus 5 (High)
The integrand is undefined at the origin, so $\lvert f\rvert$ is a continuous non-negative function
on the open set $\mathbb{R}^2\setminus\{(0,0)\}$ and \cref{def:improper_integral} applies there.
The two boundary circles are Jordan null sets, so it makes no difference whether $A$ and $B$ are
read as closed or open. Throughout, write $r := \sqrt{x^2+y^2}$, so that $\lvert f\rvert =
\lvert\cos(xy)\rvert/r^4$, and recall from \cref{rem:polar_as_shells} that a radial bound
integrates as $2\pi\int g(r)\,r\,dr$.

Nothing here can be computed exactly, because $\cos(xy)$ has no antiderivative in polar
coordinates. Both verdicts come from bounding that numerator, once from above and once from below.
\begin{enumerate}[label=\textbf{(\alph*)}]
  \item \textbf{True.} On all of $\mathbb{R}^2$, $\lvert\cos(xy)\rvert \leq 1$, so on $A$
    \[\int_A \lvert f\rvert \;\leq\; \int_A \frac{1}{r^4}
      \;=\; 2\pi\int_1^\infty \frac{1}{r^4}\,r\,dr
      \;=\; 2\pi\int_1^\infty r^{-3}\,dr
      \;=\; 2\pi\left[\frac{-1}{2r^2}\right]_1^\infty = \pi \;<\; \infty .\]
  \item \textbf{False.} Here the numerator must be bounded \emph{below}, and on $B$ it is bounded
    below by a positive constant. Indeed $\lvert xy\rvert \leq \tfrac12(x^2+y^2) \leq \tfrac12$ for
    $(x,y) \in B$ by the inequality of arithmetic and geometric means, and $\cos$ is positive and
    decreasing on $[0,\tfrac12] \subseteq [0,\pi/2)$, so $\cos(xy) \geq \cos\tfrac12 > 0$
    throughout $B$. Consequently
    \[\int_B \lvert f\rvert \;\geq\; \cos\tfrac12 \int_B \frac{1}{r^4}
      \;=\; 2\pi\cos\tfrac12 \int_0^1 r^{-3}\,dr \;=\; +\infty,\]
    since $\int_0^1 r^p\,dr$ diverges for $p \leq -1$ and here $p = -3$.
  \item \textbf{False.} $B \setminus\{(0,0)\}$ is contained in $\mathbb{R}^2\setminus\{(0,0)\}$, so
    monotonicity of the improper integral in the domain gives
    $\int_{\mathbb{R}^2}\lvert f\rvert \geq \int_B \lvert f\rvert = +\infty$ by \textbf{(b)}. One
    divergent piece is enough; the convergence found in \textbf{(a)} cannot repair it.
\end{enumerate}
\begin{remark}[The singularity is at the origin, not at infinity]
Set against \cref{ex:examHS24_TF_improper}, which asks the same two questions for
$r^{-2\alpha}$, this block is the case $\alpha = 2$: the radial integrand is $r^{-4}\cdot r =
r^{-3}$, which converges at infinity and diverges at $0$, and no exponent does both. That is the
general picture in the plane. On $A$ convergence needs $r^{-2\alpha+1}$ integrable at infinity,
that is $\alpha > 1$; on $B$ it needs the same expression integrable at $0$, that is $\alpha < 1$.
The two conditions are complementary, so $\int_{\mathbb{R}^2}r^{-2\alpha}$ diverges for every
$\alpha$, and the bounded oscillating factor $\cos(xy)$ changes none of it.
\end{remark}
\end{exercisesolution}
